\subsection{Algebraic Projections and Resolutions}
Throughout this subsection, all vector spaces are over the same fixed field $\mathbb{F}$. Projections encode direct-sum coordinates as operators, while resolutions package finite decompositions into mutually annihilating idempotents.

\begin{definition} \label{def:projection-onto-subspace-along-complement} \hlt{Projection Along a Complement}
\begin{enumerate}[label=(\roman*)]
\setlength{\itemsep}{0pt}
\item An operator $P\in\LL(\mathcal{V})$ is an \hlt{algebraic projection}, or \hlt{idempotent operator}, iff $P^2=P$.
\item If $\mathcal{V}=\mathcal{U}\oplus\mathcal{W}$, the \hlt{projection onto $\mathcal{U}$ along $\mathcal{W}$} is the map $\rho_{\mathcal{U},\mathcal{W}}:\mathcal{V}\to\mathcal{V}$ defined by $(\forall\,(\mathbf{u},\mathbf{w})\in\mathcal{U}\times\mathcal{W})\,\rho_{\mathcal{U},\mathcal{W}}(\mathbf{u}+\mathbf{w})\defeq\mathbf{u}$.
\end{enumerate}
\end{definition}

\begin{proposition} \label{prop:algebraic-projection-range-kernel-idempotent} \hlt{Algebraic Projections and Direct-Sum Splittings}\\
The following hold.
\begin{enumerate}[label=(\roman*)]
\setlength{\itemsep}{0pt}
\item For $P\in\LL(\mathcal{V})$, $P^2=P$ iff $\mathcal{V}=\operatorname{range}(P)\oplus\ker(P)$ and $P=\rho_{\operatorname{range}(P),\ker(P)}$.
\item If $\mathcal{V}=\mathcal{U}\oplus\mathcal{W}$, then $\rho_{\mathcal{U},\mathcal{W}}\in\LL(\mathcal{V})$ and $\rho_{\mathcal{U},\mathcal{W}}^2=\rho_{\mathcal{U},\mathcal{W}}$.
\item If $\mathcal{V}=\mathcal{U}\oplus\mathcal{W}$, then $\operatorname{range}(\rho_{\mathcal{U},\mathcal{W}})=\mathcal{U}$ and $\ker(\rho_{\mathcal{U},\mathcal{W}})=\mathcal{W}$.
\end{enumerate}
\begin{proof}
\begin{enumerate}[label=(\roman*)]
\setlength{\itemsep}{0pt}
\item[$(i)$] ($\Rightarrow$) Suppose $P^2=P$. For $\mathbf{v}\in\mathcal{V}$, $\mathbf{v}=P\mathbf{v}+(\mathbf{v}-P\mathbf{v})$, where $P\mathbf{v}\in\operatorname{range}(P)$ and $\mathbf{v}-P\mathbf{v}\in\ker(P)$.\\
If $\mathbf{x}\in\operatorname{range}(P)\cap\ker(P)$, choose $\mathbf{y}\in\mathcal{V}$ with $\mathbf{x}=P\mathbf{y}$. Then $\mathbf{x}=P^2\mathbf{y}=P\mathbf{x}=\mathbf{0}$, so the sum is direct by Proposition \ref{prop:uniqueness-direct-sum-two-subspaces}.\\
The displayed decomposition satisfies $(\forall\,\mathbf{v}\in\mathcal{V})\,\rho_{\operatorname{range}(P),\ker(P)}\mathbf{v}=P\mathbf{v}$; hence $P=\rho_{\operatorname{range}(P),\ker(P)}$.
\item[$(i)$] ($\Leftarrow$) Suppose $\mathcal{V}=\operatorname{range}(P)\oplus\ker(P)$ and $P=\rho_{\operatorname{range}(P),\ker(P)}$.\\
If $\mathbf{v}=\mathbf{x}+\mathbf{y}$ with $\mathbf{x}\in\operatorname{range}(P)$ and $\mathbf{y}\in\ker(P)$, then $P^2\mathbf{v}=P\mathbf{x}=\mathbf{x}=P\mathbf{v}$; hence $P^2=P$.
\item[$(ii)$] Set $R\defeq\rho_{\mathcal{U},\mathcal{W}}$ and write $\mathbf{v}_j=\mathbf{u}_j+\mathbf{w}_j$ with $(\mathbf{u}_j,\mathbf{w}_j)\in\mathcal{U}\times\mathcal{W}$ for $j=1,2$.\\
For $a,b\in\mathbb{F}$, $R(a\mathbf{v}_1+b\mathbf{v}_2)=a\mathbf{u}_1+b\mathbf{u}_2=aR\mathbf{v}_1+bR\mathbf{v}_2$ by uniqueness of direct-sum decompositions, so $R$ is linear.\\
Also $R^2(\mathbf{u}+\mathbf{w})=R\mathbf{u}=\mathbf{u}=R(\mathbf{u}+\mathbf{w})$, so $R$ is idempotent.
\item[$(iii)$] The inclusion $\operatorname{range}(\rho_{\mathcal{U},\mathcal{W}})\subseteq\mathcal{U}$ holds by the defining formula, while $\rho_{\mathcal{U},\mathcal{W}}\mathbf{u}=\mathbf{u}$ for $\mathbf{u}\in\mathcal{U}$; hence the range is $\mathcal{U}$.\\
For $\mathbf{v}=\mathbf{u}+\mathbf{w}$, $\rho_{\mathcal{U},\mathcal{W}}\mathbf{v}=\mathbf{0}$ iff $\mathbf{u}=\mathbf{0}$ iff $\mathbf{v}\in\mathcal{W}$, so the kernel is $\mathcal{W}$.
\end{enumerate}
\hlt{Structural Commentary}: Decomposition handles characterization; coordinates verify linearity, idempotence, range, and kernel.\\
\hlt{Key Insight}: Every idempotent keeps its range and kills its kernel.
\end{proof}
\end{proposition}

\begin{proposition} \label{prop:complementary-algebraic-projections} \hlt{Complementary Algebraic Projections}\\
If $\mathcal{V}=\mathcal{U}\oplus\mathcal{W}$, then the following hold.
\begin{enumerate}[label=(\roman*)]
\setlength{\itemsep}{0pt}
\item $\rho_{\mathcal{U},\mathcal{W}}+\rho_{\mathcal{W},\mathcal{U}}=I_{\mathcal{V}}$.
\item $\rho_{\mathcal{U},\mathcal{W}}\rho_{\mathcal{W},\mathcal{U}}=\rho_{\mathcal{W},\mathcal{U}}\rho_{\mathcal{U},\mathcal{W}}=0$.
\end{enumerate}
\begin{proof}
Let $\mathbf{v}=\mathbf{u}+\mathbf{w}$ with $(\mathbf{u},\mathbf{w})\in\mathcal{U}\times\mathcal{W}$.
\begin{enumerate}[label=(\roman*)]
\setlength{\itemsep}{0pt}
\item $(\rho_{\mathcal{U},\mathcal{W}}+\rho_{\mathcal{W},\mathcal{U}})\mathbf{v}=\mathbf{u}+\mathbf{w}=\mathbf{v}$. Since this equality holds on $\mathcal{V}$, the sum is $I_{\mathcal{V}}$.
\item $\rho_{\mathcal{U},\mathcal{W}}\rho_{\mathcal{W},\mathcal{U}}\mathbf{v}=\rho_{\mathcal{U},\mathcal{W}}\mathbf{w}=\mathbf{0}$ and $\rho_{\mathcal{W},\mathcal{U}}\rho_{\mathcal{U},\mathcal{W}}\mathbf{v}=\rho_{\mathcal{W},\mathcal{U}}\mathbf{u}=\mathbf{0}$. Since both equalities hold on $\mathcal{V}$, both products are zero.
\end{enumerate}
\hlt{Structural Commentary}: The two coordinate projections are evaluated on the same unique decomposition.\\
\hlt{Key Insight}: Complementary projections partition identity and annihilate each other.
\end{proof}
\end{proposition}

\begin{definition} \label{def:algebraically-orthogonal-projections-and-resolutions} \hlt{Algebraic Orthogonality and Resolutions}
\begin{enumerate}[label=(\roman*)]
\setlength{\itemsep}{0pt}
\item Algebraic projections $\rho,\sigma\in\LL(\mathcal{V})$ are \hlt{algebraically orthogonal}, written $\rho\perp\sigma$, iff $\rho\sigma=\sigma\rho=0$.
\item A \hlt{resolution of identity} on $\mathcal{V}$ is a finite family of algebraic projections $(\rho_i)_{i=1}^{k}$ such that $k\in\N$,\\
$(\forall\,i,j\in\{1,\ldots,k\})(i\neq j\Rightarrow\rho_i\perp\rho_j)$, and $\sum\limits_{i=1}^{k}\rho_i=I_{\mathcal{V}}$.
\end{enumerate}
\end{definition}

\begin{theorem} \label{thm:resolution-of-identity-gives-direct-sum} \hlt{Resolution Ranges Form a Direct Sum}\\
Let $(\rho_i)_{i=1}^{k}$ be a resolution of identity on $\mathcal{V}$, and set $\mathcal{R}_i\defeq\operatorname{range}(\rho_i)$ for $i\in\{1,\ldots,k\}$. Then the following hold.
\begin{enumerate}[label=(\roman*)]
\setlength{\itemsep}{0pt}
\item $\mathcal{V}=\bigoplus\limits_{i=1}^{k}\mathcal{R}_i$.
\item $(\forall\,i\in\{1,\ldots,k\})\,\ker(\rho_i)=\bigoplus\limits_{j\in\{1,\ldots,k\}\setminus\{i\}}\mathcal{R}_j$.
\end{enumerate}
\begin{proof}
\begin{enumerate}[label=(\roman*)]
\setlength{\itemsep}{0pt}
\item For $\mathbf{v}\in\mathcal{V}$, $\mathbf{v}=\sum\limits_{i=1}^{k}\rho_i\mathbf{v}\in\sum\limits_{i=1}^{k}\mathcal{R}_i$. Suppose $\sum\limits_{i=1}^{k}\mathbf{x}_i=\mathbf{0}$ with $\mathbf{x}_i\in\mathcal{R}_i$.\\
For $m\in\{1,\ldots,k\}$, idempotence and algebraic orthogonality imply $\rho_m\mathbf{x}_m=\mathbf{x}_m$ and $(\forall\,i\in\{1,\ldots,k\}\setminus\{m\})\,\rho_m\mathbf{x}_i=\mathbf{0}$.\\
Hence $\mathbf{x}_m=\rho_m\sum\limits_{i=1}^{k}\mathbf{x}_i=\mathbf{0}$. The sum is direct by Proposition \ref{prop:internal-direct-sum-family-criterion}.
\item By part (i), write $\mathbf{v}=\sum\limits_{j=1}^{k}\mathbf{x}_j$ uniquely with $\mathbf{x}_j\in\mathcal{R}_j$. The equality $\rho_i\mathbf{v}=\mathbf{x}_i$ holds by idempotence and algebraic orthogonality.\\
Hence $\mathbf{v}\in\ker(\rho_i)$ iff $\mathbf{x}_i=\mathbf{0}$. These are exactly the vectors in $\bigoplus\limits_{j\in\{1,\ldots,k\}\setminus\{i\}}\mathcal{R}_j$.
\end{enumerate}
\hlt{Structural Commentary}: The identity establishes spanning; orthogonal idempotents isolate each component and identify every kernel.\\
\hlt{Key Insight}: Mutually annihilating projections recover unique direct-sum coordinates.
\end{proof}
\end{theorem}

\begin{theorem} \label{thm:direct-sum-gives-resolution-of-identity} \hlt{Direct Sums Give Resolutions}\\
Let $k\in\N$, set $K\defeq\{1,\ldots,k\}$, and suppose $\mathcal{V}=\bigoplus\limits_{i\in K}\mathcal{U}_i$. For $i\in K$, set $\rho_i\defeq\rho_{\mathcal{U}_i,\sum\limits_{j\in K\setminus\{i\}}\mathcal{U}_j}$. Then $(\rho_i)_{i\in K}$ is a resolution of identity.
\begin{proof}
For $i\in K$, set $\mathcal{W}_i\defeq\sum\limits_{j\in K\setminus\{i\}}\mathcal{U}_j$. Then $\mathcal{U}_i\cap\mathcal{W}_i=\{\mathbf{0}\}$ by Proposition \ref{prop:internal-direct-sum-family-criterion}(iii), while $\mathcal{V}=\mathcal{U}_i+\mathcal{W}_i$ by the assumed sum.\\
Hence $\mathcal{V}=\mathcal{U}_i\oplus\mathcal{W}_i$ by Proposition \ref{prop:uniqueness-direct-sum-two-subspaces}, and $\rho_i$ is an algebraic projection by Proposition \ref{prop:algebraic-projection-range-kernel-idempotent}(ii).\\
For the unique decomposition $\mathbf{v}=\sum\limits_{i\in K}\mathbf{u}_i$, $\rho_i\mathbf{v}=\mathbf{u}_i$; hence $(\sum\limits_{i\in K}\rho_i)\mathbf{v}=\mathbf{v}$. This pointwise identity proves $\sum\limits_{i\in K}\rho_i=I_{\mathcal{V}}$.\\
If $i\neq j$, then $\rho_i\rho_j\mathbf{v}=\rho_i\mathbf{u}_j=\mathbf{0}$ and $\rho_j\rho_i\mathbf{v}=\rho_j\mathbf{u}_i=\mathbf{0}$. Thus $\rho_i\perp\rho_j$, and the family is a resolution of identity.\\
\hlt{Structural Commentary}: Direct-sum uniqueness defines coordinate projections; pointwise identities verify the resolution conditions.\\
\hlt{Key Insight}: Coordinate projections isolate summands and reconstruct every vector.
\end{proof}
\end{theorem}

\begin{theorem} \label{thm:algebra-of-projections} \hlt{Algebra of Projections}\\
Let $\rho,\sigma\in\LL(\mathcal{V})$ be algebraic projections, and assume $2\neq0$ in $\mathbb{F}$.
\begin{enumerate}[label=(\roman*)]
\setlength{\itemsep}{0pt}
\item The operator $\rho+\sigma$ is an algebraic projection iff $\rho\perp\sigma$.\\
Under these conditions, $\operatorname{range}(\rho+\sigma)=\operatorname{range}(\rho)\oplus\operatorname{range}(\sigma)$ and $\ker(\rho+\sigma)=\ker(\rho)\cap\ker(\sigma)$.
\item The operator $\rho-\sigma$ is an algebraic projection iff $\rho\sigma=\sigma\rho=\sigma$.\\
Under these conditions, $\operatorname{range}(\rho-\sigma)=\operatorname{range}(\rho)\cap\ker(\sigma)$ and $\ker(\rho-\sigma)=\ker(\rho)\oplus\operatorname{range}(\sigma)$.
\item If $\rho\sigma=\sigma\rho$, then $\rho\sigma$ is an algebraic projection, $\operatorname{range}(\rho\sigma)=\operatorname{range}(\rho)\cap\operatorname{range}(\sigma)$, and $\ker(\rho\sigma)=\ker(\rho)+\ker(\sigma)$.
\end{enumerate}
\begin{proof}
\begin{enumerate}[label=(\roman*)]
\setlength{\itemsep}{0pt}
\item[$(i)$] ($\Rightarrow$) Set $A\defeq\rho\sigma$, $B\defeq\sigma\rho$, and $C\defeq\rho\sigma\rho$. The identity $(\rho+\sigma)^2=\rho+\sigma$ is equivalent to $A+B=0$.\\
Then $A+C=\rho(A+B)=0$ and $C+B=(A+B)\rho=0$, so $A=B$ and $2A=0$.\\
The scalar $2$ is invertible by $2\neq0$ in the field $\mathbb{F}$; hence $A=B=0$ and $\rho\perp\sigma$.
\item[$(i)$] ($\Leftarrow$) Suppose $\rho\perp\sigma$. Then $(\rho+\sigma)^2=\rho+\sigma$, so the sum is an algebraic projection.\\
Its range is contained in $\operatorname{range}(\rho)+\operatorname{range}(\sigma)$, while $(\rho+\sigma)(\mathbf{x}+\mathbf{y})=\mathbf{x}+\mathbf{y}$ for $(\mathbf{x},\mathbf{y})\in\operatorname{range}(\rho)\times\operatorname{range}(\sigma)$.\\
If $\mathbf{z}\in\operatorname{range}(\rho)\cap\operatorname{range}(\sigma)$ and $\mathbf{z}=\sigma\mathbf{w}$, then $\mathbf{z}=\rho\mathbf{z}=\rho\sigma\mathbf{w}=\mathbf{0}$; hence the range sum is direct by Proposition \ref{prop:uniqueness-direct-sum-two-subspaces}.\\
$(\rho+\sigma)\mathbf{v}=\mathbf{0}$ $\Rightarrow$ $\rho\mathbf{v}=\rho(\rho+\sigma)\mathbf{v}=\mathbf{0}$ and $\sigma\mathbf{v}=\sigma(\rho+\sigma)\mathbf{v}=\mathbf{0}$. Reverse implication is immediate. Thus $\ker(\rho+\sigma)=\ker(\rho)\cap\ker(\sigma)$.
\item[$(ii)$] ($\Rightarrow$) Suppose $Q\defeq\rho-\sigma$ is algebraic projection. Then $(I_{\mathcal{V}}-Q)^2=I_{\mathcal{V}}-Q$ and $(I_{\mathcal{V}}-\rho)^2=I_{\mathcal{V}}-\rho$ hold by $Q^2=Q$ and $\rho^2=\rho$.\\
Since $I_{\mathcal{V}}-Q=(I_{\mathcal{V}}-\rho)+\sigma$, $(I_{\mathcal{V}}-\rho)\perp\sigma$ by part (i). Therefore $(I_{\mathcal{V}}-\rho)\sigma=\sigma(I_{\mathcal{V}}-\rho)=0$, which is equivalent to $\rho\sigma=\sigma\rho=\sigma$.
\item[$(ii)$] ($\Leftarrow$) Suppose $\rho\sigma=\sigma\rho=\sigma$ and set $Q\defeq\rho-\sigma$. The operator $I_{\mathcal{V}}-\rho$ is an algebraic projection and is algebraically orthogonal to $\sigma$, so $I_{\mathcal{V}}-Q=(I_{\mathcal{V}}-\rho)+\sigma$ is an algebraic projection by part (i). Since $I_{\mathcal{V}}-Q$ is idempotent, $Q^2=(I_{\mathcal{V}}-(I_{\mathcal{V}}-Q))^2=I_{\mathcal{V}}-(I_{\mathcal{V}}-Q)=Q$.\\
Since $\rho Q=Q$ and $\sigma Q=0$, $\operatorname{range}(Q)\subseteq\operatorname{range}(\rho)\cap\ker(\sigma)$; the reverse inclusion holds because $Q\mathbf{x}=\mathbf{x}$ when $\rho\mathbf{x}=\mathbf{x}$ and $\sigma\mathbf{x}=\mathbf{0}$.\\
The operator $Q$ kills both $\ker(\rho)$ and $\operatorname{range}(\sigma)$. Their intersection is trivial: if $\mathbf{z}=\sigma\mathbf{w}\in\ker(\rho)$, then $\mathbf{z}=\rho\sigma\mathbf{w}=\rho\mathbf{z}=\mathbf{0}$.\\
Their sum is direct by Proposition \ref{prop:uniqueness-direct-sum-two-subspaces}, and $\ker(\rho)\oplus\operatorname{range}(\sigma)\subseteq\ker(Q)$ because $Q$ kills both summands.\\
Conversely, if $\mathbf{v}\in\ker(Q)$, then $\rho\mathbf{v}=\sigma\mathbf{v}$ and $\mathbf{v}=(\mathbf{v}-\rho\mathbf{v})+\sigma\mathbf{v}\in\ker(\rho)\oplus\operatorname{range}(\sigma)$. Hence $\ker(Q)=\ker(\rho)\oplus\operatorname{range}(\sigma)$.
\item[$(iii)$] Suppose $\rho\sigma=\sigma\rho$. Then $(\rho\sigma)^2=\rho^2\sigma^2=\rho\sigma$, so $\rho\sigma$ is an algebraic projection.\\
If $\mathbf{x}=\rho\sigma\mathbf{y}$, then $\rho\mathbf{x}=\mathbf{x}=\sigma\mathbf{x}$; conversely, $\rho\mathbf{x}=\mathbf{x}=\sigma\mathbf{x}$ implies $\mathbf{x}=\rho\sigma\mathbf{x}$. Thus $\operatorname{range}(\rho\sigma)=\operatorname{range}(\rho)\cap\operatorname{range}(\sigma)$.\\
If $\mathbf{x}\in\ker(\rho)$, then $\rho\sigma\mathbf{x}=\sigma\rho\mathbf{x}=\mathbf{0}$; if $\mathbf{x}\in\ker(\sigma)$, then $\rho\sigma\mathbf{x}=\mathbf{0}$. Hence $\ker(\rho)+\ker(\sigma)\subseteq\ker(\rho\sigma)$.\\
If $\rho\sigma\mathbf{v}=\mathbf{0}$, then $\mathbf{v}=(\mathbf{v}-\sigma\mathbf{v})+\sigma\mathbf{v}$ with $\mathbf{v}-\sigma\mathbf{v}\in\ker(\sigma)$ and $\sigma\mathbf{v}\in\ker(\rho)$; hence $\ker(\rho\sigma)=\ker(\rho)+\ker(\sigma)$.
\end{enumerate}
\hlt{Structural Commentary}: Expanding idempotence isolates cross terms; fixed and killed subspaces determine each range and kernel.\\
\hlt{Key Insight}: Commuting idempotents encode intersections and sums.
\end{proof}
\end{theorem}
